\documentclass[12pt]{article}
\usepackage{amsmath, amsfonts, amsthm}
\usepackage[margin=1in]{geometry}
\newcommand{\Mod}[1]{\ (\mathrm{mod}\ #1)}
\newcommand{\field}[1]{\mathbb{#1}}
\newcommand{\Q}{\field{Q}}
\newcommand{\R}{\field{R}}
\newcommand{\N}{\field{N}}
\newcommand{\C}{\field{C}}
\newcommand{\Z}{\field{Z}}
\newcommand{\F}{\field{F}}
\title{Session 2 Solutions: Solving polynomial equations modulo a prime}
\date{February 15, 2026}

\begin{document}
\maketitle

\begin{enumerate}

%% Problem 1
\item
\begin{enumerate}
  \item[(a)] By Fermat's Little Theorem, $3^{16} \equiv 1 \Mod{17}$.
  Write $2027 = 16 \cdot 126 + 11$, so $3^{2027} \equiv 3^{11} \Mod{17}$.
  Now $3^2=9$, $3^4=81\equiv 81-4\cdot17=13\equiv -4$, $3^8\equiv 16\equiv -1$,
  so $3^{11}=3^8\cdot 3^2\cdot 3 \equiv (-1)(9)(3)=-27\equiv -27+2\cdot17=7\Mod{17}$.

  \textbf{Answer: }$\mathbf{7}$.

  \item[(b)] By FLT, $7^{22}\equiv 1\Mod{23}$.
  Write $1958 = 22\cdot 89 + 0$, so $7^{1958}\equiv 7^0 \equiv 1\Mod{23}$.

  \textbf{Answer: }$\mathbf{1}$.

  \item[(c)] By FLT, $4^{30}\equiv 1\Mod{31}$.
  Write $691 = 30\cdot 23 + 1$, so $4^{691}\equiv 4^1 \equiv 4\Mod{31}$.

  \textbf{Answer: }$\mathbf{4}$.
\end{enumerate}

%% Problem 2
\item The last digit of $7^{100}$ is $7^{100}\Mod{10}$.
The powers of $7$ mod $10$ cycle with period $4$: $7,9,3,1,7,\ldots$
Since $100=4\cdot 25$, we get $7^{100}\equiv 1\Mod{10}$.

\textbf{Answer: }$\mathbf{1}$.

%% Problem 3
\item We must show $42 \mid n^7-n$. Note $42=2\cdot 3\cdot 7$, so it suffices to show $2$, $3$, and $7$ each divide $n^7-n$.

By Fermat's Little Theorem, for any prime $p$, $n^p \equiv n \Mod{p}$, i.e.\ $p\mid n^p - n$.

Taking $p=7$: $7\mid n^7-n$.

For $p=3$: $3\mid n^3-n$. Now $n^7-n = n(n^6-1)=n(n^3-1)(n^3+1)$.
Since $n^3\equiv n\Mod{3}$, we get $n^3-1\equiv n-1$ and $n^3+1\equiv n+1\Mod{3}$,
so $n^7-n\equiv n(n-1)(n+1)=n^3-n\equiv 0\Mod{3}$.

For $p=2$: similarly $2\mid n^2-n$ and $n^7-n = n(n^6-1)$. Since $n^2\equiv n\Mod{2}$, $n^6\equiv n\Mod{2}$, so $2\mid n^7-n$.

Since $2$, $3$, and $7$ each divide $n^7-n$ and $2\cdot 3\cdot 7=42$, we get $42\mid n^7-n$. $\qed$

%% Problem 4
\item
\begin{enumerate}
  \item[(a)] We need $5x\equiv 6\Mod{8}$. We look for the inverse of $5$ modulo $8$:
  $5\cdot 5=25\equiv 1\Mod{8}$, so $5^{-1}\equiv 5$.
  Then $x\equiv 5\cdot 6=30\equiv 6\Mod{8}$.

  \textbf{Answer: }$\mathbf{x\equiv 6\Mod{8}}$.

  \item[(b)] We need $12x\equiv 5\Mod{19}$. We find $12^{-1}\Mod{19}$.
  By Fermat's Little Theorem, $12^{-1}\equiv 12^{17}\Mod{19}$.
  Computing: $12\equiv -7$, $(-7)^2=49\equiv 11$, $11^2=121\equiv 121-6\cdot 19=7$, $7^2=49\equiv 11$, so $12^8\equiv 11$.
  Then $12^{16}\equiv 11^2=121\equiv 7$ and $12^{17}\equiv 7\cdot 12=84\equiv 84-4\cdot 19=8$.
  So $12^{-1}\equiv 8\Mod{19}$ and $x\equiv 8\cdot 5=40\equiv 2\Mod{19}$.

  \textbf{Answer: }$\mathbf{x\equiv 2\Mod{19}}$.

  \item[(c)] We find all $(x,y)$ with $x^2+y^2\equiv 1\Mod{13}$.
  The quadratic residues mod $13$ are $\{0,1,3,4,9,10,12\}$ (i.e.\ $0^2=0$, $1^2=1$, $2^2=4$, $3^2=9$, $4^2=3$, $5^2=12$, $6^2=10$).

  For each $x\in\{0,\ldots,12\}$, we need $y^2\equiv 1-x^2\Mod{13}$.
  This is solvable iff $1-x^2$ is a quadratic residue mod $13$.

  One can check all cases: there are $\mathbf{12}$ solutions $(x,y)$ in $(\Z/13\Z)^2$.
  For example, $(x,y)=(0,1),(0,12),(1,0),(12,0),(2,6),(2,7),(6,2),(6,11),(7,2),(7,11),(11,6),(11,7)$.
\end{enumerate}

%% Problem 5
\item Suppose $x^3+y^3=z^3$ for integers $x,y,z$.
The cubes modulo $7$ are: $0^3\equiv 0$, $1^3\equiv 1$, $2^3\equiv 1$, $3^3\equiv 6$, $4^3\equiv 1$, $5^3\equiv 6$, $6^3\equiv 6\Mod{7}$.
So the only possible values of $n^3\Mod{7}$ are $\{0,1,6\}$.

If none of $x,y,z$ is divisible by $7$, then $x^3,y^3,z^3\in\{1,6\}\Mod{7}$.
We check all possibilities of $a+b\equiv c\Mod{7}$ where $a,b,c\in\{1,6\}$:
$1+1=2$, $1+6=0$, $6+1=0$, $6+6=5$. None of $\{2,0,0,5\}$ lie in $\{1,6\}$ ---
except $0$, but $0\notin\{1,6\}$. So no combination works, a contradiction.

Hence at least one of $x,y,z$ must be divisible by $7$. $\qed$

%% Problem 6
\item
\begin{enumerate}
  \item[(a)] $(x+1)^5 = x^5 + 5x^4 + 10x^3 + 10x^2 + 5x + 1$.

  \item[(b)] By the binomial theorem, $(x+y)^p = \sum_{k=0}^{p}\binom{p}{k}x^k y^{p-k}$.
  For $0<k<p$, $\binom{p}{k}=\frac{p!}{k!(p-k)!}$. Since $p$ is prime and $p$ divides the numerator but not the denominator ($1\le k\le p-1$), we get $p\mid\binom{p}{k}$.
  So modulo $p$, all middle terms vanish and $(x+y)^p\equiv x^p+y^p\Mod{p}$. $\qed$
\end{enumerate}

%% Problem 7
\item We show $3x^3+4y^3\equiv 1715\Mod{19}$ has a solution.
Note $1715 = 19\cdot 90 + 5$, so we need $3x^3+4y^3\equiv 5\Mod{19}$.
Try $x=0$: $4y^3\equiv 5\Mod{19}$. We need $4^{-1}\Mod{19}$: $4\cdot 5=20\equiv 1$, so $4^{-1}\equiv 5$.
Then $y^3\equiv 25\equiv 6\Mod{19}$. Check: $3^3=27\equiv 27-19=8$. $4^3=64\equiv 64-3\cdot 19=7$. $10^3=1000\equiv 1000-52\cdot 19=1000-988=12$. $11^3=1331\equiv 1331-70\cdot 19=1331-1330=1$. $15^3=3375\equiv 3375-177\cdot 19=3375-3363=12$. $8^3=512\equiv 512-26\cdot 19=512-494=18\equiv -1$. $16^3\equiv(-3)^3=-27\equiv -8\equiv 11$. $5^3=125\equiv 125-6\cdot 19=125-114=11$. $6^3=216\equiv 216-11\cdot 19=216-209=7$. $13^3\equiv(-6)^3=-7\equiv 12$. $9^3=729\equiv 729-38\cdot 19=729-722=7$.

Let us try $(x,y)=(1,y)$: $3+4y^3\equiv 5$, so $4y^3\equiv 2$, $y^3\equiv 5\cdot 2=10\Mod{19}$. From the cube table: $12^3\equiv(-7)^3=-343\equiv -343+19\cdot 19=-343+361=18$. $7^3=343\equiv 343-18\cdot 19=343-342=1$.

Try $(x,y)=(3,y)$: $3\cdot 27+4y^3=81+4y^3\equiv 81-4\cdot 19+4y^3=5+4y^3\equiv 5$, so $4y^3\equiv 0$, giving $y\equiv 0\Mod{19}$.

\textbf{Answer: }$(x,y)\equiv (3,0)\Mod{19}$ is a solution, since $3\cdot 27+4\cdot 0=81=4\cdot 19+5\equiv 5\equiv 1715\Mod{19}$.

%% Problem 8
\item
\begin{enumerate}
  \item[(a)] OR$(x,y)$ = NOT(AND(NOT($x$), NOT($y$))) by De Morgan's law.
  NOT($x$) $=1+x$, NOT($y$)$=1+y$, AND $=(1+x)(1+y)=1+x+y+xy$.
  NOT of that $=1+(1+x+y+xy)=x+y+xy$.

  (Over $\F_2$: check $P(0,0)=0$, $P(1,0)=1$, $P(0,1)=1$, $P(1,1)=1+1+1=1$. Correct.)

  \textbf{Answer: }$P(x,y)=x+y+xy$.

  \item[(b)] ``($x$ AND $y$) is True'' becomes $xy=1$.
  ``($x$ OR $z$) is False'' means OR$(x,z)=0$, i.e.\ $x+z+xz=0$.

  \textbf{System:} $xy=1$ and $x+z+xz=0$ over $\F_2$.
\end{enumerate}

%% Problem 9
\item
\begin{enumerate}
  \item We check each $x\in\{0,1,\ldots,6\}$:

  \begin{tabular}{c|c|c|c}
  $x$ & $x^3+2\Mod{7}$ & Is it a square? & $y$ \\ \hline
  0 & 2 & Yes ($3^2\equiv 2$) & $3,4$ \\
  1 & 3 & No & --- \\
  2 & 3 & No & --- \\
  3 & 1 & Yes & $1,6$ \\
  4 & 3 & No & --- \\
  5 & 1 & Yes & $1,6$ \\
  6 & 1 & Yes & $1,6$
  \end{tabular}

  (Squares mod $7$: $0^2=0$, $1^2=1$, $2^2=4$, $3^2=2$, $4^2=2$, $5^2=4$, $6^2=1$. So the squares in $\F_7$ are $\{0,1,2,4\}$.)

  The affine points are: $(0,3),(0,4),(3,1),(3,6),(5,1),(5,6),(6,1),(6,6)$.

  \textbf{8 affine points.}

  \item Including the point at infinity $\mathcal{O}$, the total is $\mathbf{9}$ points.
\end{enumerate}

\end{enumerate}
\end{document}
